\section{Теоретическая часть и математические основы}

Рассмотрим базовые математические закономерности, используемые при анализе данных. Первой важной формулой является фундаментальное тождество Эйлера, связывающее пять ключевых математических констант:
\begin{equation}
\label{eq:euler}
e^{i\pi} + 1 = 0
\end{equation}

Далее определим функцию нормального распределения вероятностей (кривая Гаусса), которая повсеместно применяется при статистическом анализе данных и в машинном обучении~\cite{henrick2017ml}. Плотность распределения имеет вид:
\begin{equation}
\label{eq:gauss}
f(x) = \frac{1}{\sigma \sqrt{2\pi}} \exp\left( -\frac{1}{2}\left(\frac{x-\mu}{\sigma}\right)^{\!2} \, \right)
\end{equation}

Для оценки итерационных процессов алгоритмов нам потребуется аппарат дискретных сумм. Запишем формулу суммы первых $n$ натуральных чисел:
\begin{equation}
\label{eq:sum}
\sum_{i=1}^{n} i = \frac{n(n+1)}{2}
\end{equation}

Все три формулы "--- \eqref{eq:euler}, \eqref{eq:gauss} и \eqref{eq:sum} "--- являются классическими примерами математического анализа и широко используются при проектировании сложных вычислительных систем.

\subsection{Сводные аналитические данные}

В таблице~\ref{tab:metrics} приведены результаты сравнительного анализа производительности алгоритмов при изменении управляющей переменной \texttt{N}.

\begin{table}[htbp]
    \centering
    \caption{Сравнение времени работы алгоритмов вычислений}
    \label{tab:metrics}
    \begin{tabular}{ccc}
        \hline
        Размерность (\texttt{N}) & Алгоритм А (сек) & Алгоритм Б (сек) \\
        \hline
        100  & 0.001 & 0.002 \\
        1000 & 0.015 & 0.008 \\
        5000 & 0.421 & 0.095 \\
        \hline
    \end{tabular}
\end{table}